% =============================================================================
%  Part II -- The causal picture of the world (ontology)
% =============================================================================

This part describes the world as Causal Physics sees it, in words rather than equations. It follows
one line of thought: from what exists, through time, space and motion, to matter, gravity and the
limits of what can be known. The mathematics of the following parts is a precise form of this
picture. Each statement carries a label that says how firmly it stands:
\textsf{[tested]} -- supported by a test of the laboratory (identifier given);
\textsf{[derived]} -- follows from a derivation in this paper or from established results it uses;
\textsf{[observed]} -- seen in a test, but not the outcome the test judged;
\textsf{[interpretation]} -- a reading consistent with the results, from which nothing is computed;
\textsf{[open]} -- part of the programme, not yet formulated or tested.

% =============================================================================
\section{What exists: events and influence}\label{sec:onto-events}
% =============================================================================

The world is not made of things that exist in space and persist in time. It is made of
\emph{events}: elementary happenings, each of which is an act of realisation. An event has no size,
no position, no mass and no charge. What it has is relations: some events can influence it, and it
can influence others. Everything we call a property, the position of a stone, the mass of a planet,
the tick of a clock, is a pattern in these relations, visible only when many events are taken
together. \stat{interpretation}

The single primitive relation is \emph{influence}: one event can affect another. Influence has a
direction, and it never closes into a loop: nothing influences its own past. With these two facts
the set of all events becomes an \emph{order}, and the order is the deepest structure of the world.
Everything else in this part is a way of reading it. \stat{derived}

A table, in this picture, is not an object that happens to be at a place. It is a pattern of events
that keeps reproducing itself, update after update. Its solidity is the stability of the pattern,
and its place is the way the pattern is woven into the surrounding web of influence.
\stat{interpretation}

% =============================================================================
\section{Time: the order seen from inside}\label{sec:onto-time}
% =============================================================================

``Earlier'' means ``could have influenced''; ``later'' means ``could have been influenced by''. Time
is nothing more than this. There is no universal now: two events that cannot influence each other
have no fact about which came first, and no observer can find one. \stat{derived}

Each observer has a local time, the number of events along their own history. A clock is a pattern
that repeats, and its ticks are counted events. The clocks of two observers agree only where their
histories are related, and they need not agree on how much time has passed between two meetings.
\stat{derived}

The passage of time is the growth of the order. New events appear in the future of existing ones and
never in their past. The past is fixed because it has already been realised; the future is open
because it has not. \stat{interpretation}

% =============================================================================
\section{Space and dimension}\label{sec:onto-space}
% =============================================================================

Space is what lies between events that cannot influence each other. A ``moment'' of the world, as
an observer experiences it, is a large set of mutually unrelated events, and the distances within it
are read from how far one has to go into the past or the future to find events related to both. Space
is therefore not a container that exists first; it is the side-by-side structure of the order.
\stat{derived}

The order fixes the shape of spacetime up to a local scale, and counting events fixes the scale.
Nothing else is needed: distances, volumes, the rate of clocks and the dimension itself are read
from order and number. The dimension, for example, is visible in what fraction of pairs of events
in a region are related. \stat{tested: K.1, K.16b, K.16c}

Why space has three dimensions is not derived here. But it is notable that the way the order responds
to matter (Sec.~\ref{sec:onto-gravity}) agrees with the gravity we observe in exactly three spatial
dimensions and in no other. \stat{derived}\quad Whether this is the reason for three dimensions is
an open question. \stat{open}

% =============================================================================
\section{Motion and relativity}\label{sec:onto-motion}
% =============================================================================

The substrate looks the same to every observer moving uniformly. There is no state of rest in it, no
ether against which motion could be measured, and no experiment can find one. This is the principle of
relativity, and here it is a property of the substrate, not an additional postulate.
\stat{derived}

The speed of influence is the same for all observers for the same reason: the boundary between
events that can and cannot influence each other is part of the order, and the order is the same for
everyone. The speed of light is the conversion between counting events along time and counting them
across space. \stat{derived}

A moving clock runs slow because it takes a shorter path through the order. Between two meetings,
the observer who stays at rest passes through the longest chain of events, and an observer who
travels away and comes back passes through fewer. The travelling twin is younger because fewer
events happened to her, not because anything acted on her clock. \stat{derived}

% =============================================================================
\section{Observing is receiving causes}\label{sec:onto-observe}
% =============================================================================

To observe something is to be on the receiving end of its influence. Seeing a tree means that
events in the tree have influenced, through a long chain of intermediate events, the events of our
eye. There is no observation without causal contact, and every observation is itself an event that
joins the order.

This explains why the world looks different near and far. A tree next to us is rich in detail; the
same tree on the horizon is a silhouette. The tree has not changed. What has changed is the density
and fidelity of the influence that connects it to us. Two different effects are at work.
\begin{itemize}
\item \emph{Geometric dilution.} The same influence spreads over a larger sphere as it travels, so
  less of it reaches us, in proportion to $1/r^2$. This is the root of the inverse-square law.
  \stat{tested: F1.1}
\item \emph{Loss through intermediaries.} Where influence passes through many intermediate events
  (air, dust, memory, a historical record), each of them reshuffles it, and detail is lost. This is
  the entropic side of distance. \stat{interpretation}
\end{itemize}
Influence that travels through empty space along the light cone does not pass through intermediaries
in this sense: events on the light cone are causally close however distant they are in space. The
images of very distant point sources are therefore not blurred by the discreteness of the substrate,
and the sharp images of distant quasars agree with that. \stat{derived}

% =============================================================================
\section{Scale and causal resolution}\label{sec:onto-scale}
% =============================================================================

What we call scale is the resolution with which the order can be read. Different physical regimes
are different limits of that resolution, not different sizes.

\paragraph{The discreteness scale.} At the scale of individual events the order is well defined, but
there are too few events in any region to fix a shape. Trajectories, fields and distances lose their
meaning here, and only the relations remain. \stat{derived}

\paragraph{The plateau.} When regions contain many events and the order is close to a uniform
sprinkling, geometry becomes an excellent compression of it. This is the regime of ordinary matter
and of everyday experience, in which three-dimensional space and Lorentzian spacetime are reliable
descriptions. Dynamics of orders of this kind has a phase whose typical members look like spacetime,
separated by a transition from phases that do not.
\stat{known in 2D~\cite{Surya2012}; in part F6.4a}\quad
General relativity and quantum field theory are effective descriptions on the plateau.
\stat{interpretation}

\paragraph{Between the two.} At an intermediate scale $\ell_k$ the order is smooth enough for
geometry but still remembers that it is discrete. Its response reaches along the light cone to
distant events, and the response to a static field depends on how long that field has existed.
\stat{derived, Sec.~\ref{sec:bdaction}}

% =============================================================================
\section{Transport: nothing from nothing}\label{sec:onto-transport}
% =============================================================================

The substrate is not a passive stage. It carries influence from events to the events they affect,
and this influence is conserved: what flows into an event flows out of it, apart from what the event
stores. Nothing appears from nothing and nothing disappears into nothing. This continuity is the
most basic law of the substrate, and the conservation laws of physics are its large-scale
expressions. A steady source surrounded by any closed surface sends exactly its own output through
it, which is Gauss's law. \stat{tested: F1.1}

The capacity of an event to pass influence on is finite. When more influence arrives than an event
can carry, it is held back: the substrate buffers it, and an observer sees this as delay. Delay is
not a defect of transmission but the visible expression of a finite causal bandwidth.
\stat{interpretation}

% =============================================================================
\section{Branching, merging and the meaning of energy}\label{sec:onto-merging}
% =============================================================================

Influence branches and merges. \emph{Branching} distributes the influence of one event among many
later events. \emph{Merging} compresses influence that arrived along distinct paths into a single
continuation. At the level of the substrate the two are symmetric: every event is a meeting point of
its past and a source of its future. \stat{interpretation}

\paragraph{Energy.} The amount of influence carried is \emph{energy}; the direction in which it is
carried is \emph{momentum}. Together they are one quantity seen from different frames: energy is the
part along time, momentum the part across space. Both are conserved because influence is.
\stat{tested: K.15}

\paragraph{Mass.} When two streams of influence moving in different directions merge, their directed
motion is not lost. It becomes internal: the merged stream moves more slowly as a whole and carries
more rest energy, in exactly the proportion that conserves energy and momentum together. Merging is
how motion becomes mass and heat; two counter-streaming populations that merge cool down. Mass is
energy that no longer points anywhere. This is the content of $E=mc^2$ in this picture.
\stat{tested: K.15, F4.1}

% =============================================================================
\section{Light and matter}\label{sec:onto-matter}
% =============================================================================

\paragraph{Light.} Influence that travels on the boundary of the order, along the light cone, has no
clock: no events happen along its way, it cannot be at rest, and it has no rest frame. Light is
influence in pure transit. \stat{derived}

\paragraph{Matter.} Matter is influence that persists. A flow that only spread out would be radiation:
it would fill its light cone and dissolve. A flow that stays together, carried from event to event
along a narrow bundle of histories, is matter, and along its history events happen: matter has
clocks.

What keeps a flow together on a substrate that has no preferred directions is \emph{memory}. A flow
that carries its own direction of motion as a state passed from event to event stays on the world line
of its source and moves with the source when the source moves. This memory of direction is
\emph{inertia}: a body continues in its motion because its flow remembers where it was going.
Momentum is not added to matter from outside; it is what makes localised matter possible at all.
\stat{tested: F3.6}

\paragraph{Structure.} Where repeated activity keeps routing influence through the same region, the
substrate approaches its capacity and influence can no longer spread freely: it is confined to a
restricted set of paths. Structure is the self-maintaining outcome of this confinement, a dynamical
equilibrium between the continuity of influence and the finiteness of capacity, maintained as long as
the activity that creates it continues. \stat{interpretation}\quad How such confinement closes into
a self-bound object, a particle that holds itself together without an external source, is the central
open problem of the matter sector. \stat{open}

% =============================================================================
\section{Gravity: the order responds to matter}\label{sec:onto-gravity}
% =============================================================================

\paragraph{Load.} Matter is a flow that passes through the substrate continuously, and the substrate
carries it as load. A loaded substrate does three things at once: its clocks run slower, because more
of its activity goes into carrying the flow; its causal cones narrow, because influence travels more
slowly through it; and it becomes denser, because it adds events.
\stat{effects derived; reasons: interpretation}

\paragraph{Throughput.} The last two effects are tied together by one principle, familiar from traffic
and from every conserved current passing through a constriction: \emph{the throughput is conserved}.
Where cars slow down they bunch up, so that the number passing a point per minute stays the same.
Where the substrate slows influence down, it adds events, so that the number of events a light front
passes per unit time and area stays the same. The speed of light is the maximal rate at which
influence travels without congestion; near matter the effective speed falls below it.
\stat{derived; tested: F7.1c}

\paragraph{The field.} The response of the order is not confined to where the matter is. It spreads,
at the speed of influence, into the surrounding substrate, weakening with distance as influence does.
The gravitational field is this spreading response of the order: not a substance added to spacetime,
but the order itself, reshaped. \stat{tested: F1.3}

\paragraph{Falling.} A body falls because its history bends towards slower clocks: of all histories
between two events, matter follows the one along which the most events happen, and near a mass that
history leans towards the mass. Light bends because the cones it travels along narrow towards
matter. The amount by which light bends, twice what a naive Newtonian picture gives, is exactly what
the conservation of throughput gives in three spatial dimensions. \stat{derived; fall from the order: 2+1 passed (F5.1d), 3+1 failed (F5.1e)}

\paragraph{Equivalence.} All bodies fall in the same way, whatever they are made of, because falling is
not a reaction to a force but the following of the order's own longest histories. What a body is made
of does not enter into which history is longest. Being pushed and being attracted probe the same
property of the substrate from two sides. \stat{derived}

\paragraph{Everything that moves gravitates.} The substrate responds to a flow in the flow's own rest
frame, and the responses of different flows add. Because of this, not only mass but also pressure
and motion gravitate, and a moving flow drags the causal cones along with it.
\stat{tested: F7.4a}\quad The order responds in this way because its own dynamics, the action that
counts its intervals, has exactly the stiffness of Einstein's gravity.
\stat{derived and computed, Sec.~\ref{sec:bdaction}}

\paragraph{Non-linearity.} Weak fields add. Strong ones do not: when two sources sit close together,
each lies in the well of the other and acts more weakly, so the common well is shallower than the sum
of the two. Gravity is non-linear because the order responds to its own response.
\stat{tested: F1.4}

% =============================================================================
\section{Horizons and the limits of reconstructability}\label{sec:onto-horizons}
% =============================================================================

Where the substrate is loaded very strongly, geometry stops being a good description of the order.
The events are still there and influence still flows, but the order no longer behaves like a smooth
spacetime: reading it with clocks, cones and densities gives answers that do not fit together into
one metric. \stat{observed: F7.1b}

A \emph{causal horizon} is the limit of reconstructability: the boundary beyond which an external
observer can no longer keep an order-preserving description of what happens inside using a global
time, a smooth geometry or isotropic propagation. Beyond it, causal updates continue without
interruption and causal continuity is intact. What is lost is the ability to describe the inside in
the language of the outside. Horizons mark the reach of geometric language; they are not a boundary
of causal evolution itself. \stat{interpretation}

Merging plays a central role here. When many causal channels converge under extreme load, the causal
distance between events collapses from the substrate's own point of view. To an observer who relies on
geometric reconstruction, this collapse appears as temporal ambiguity, loss of detail and horizon
behaviour. \stat{interpretation}

% =============================================================================
\section{Entropy and the arrow of time}\label{sec:onto-entropy}
% =============================================================================

The order has a direction because it grows: events are added to its future, never to its past. This
is the deepest arrow of time. \stat{interpretation}

Merging gives the arrow a second, familiar face. Branching can be followed backwards: every branch
remembers where it came from. Merging cannot be undone uniquely: once distinct paths have joined
into one continuation, many different pasts are compatible with the same present. What is lost is
the record of which past it was, and the measure of that loss is \emph{entropy}. Irreversibility is
the asymmetry between how easy it is to merge and how hard it is to unmerge. \stat{interpretation}

A horizon carries an entropy proportional to its area, and the order counts it directly: the number
of pairs of events that straddle the horizon in a minimal way is proportional to its proper area.
\stat{tested: K.14}\quad This is the bridge to the thermodynamic view of gravity, in which Einstein's
equations describe the equilibrium of the order the way the gas laws describe the equilibrium of
molecules. \stat{open, Sec.~\ref{sec:jacobson}}

% =============================================================================
\section{The vacuum is not empty}\label{sec:onto-vacuum}
% =============================================================================

Empty space is not nothing. It is the substrate with no flow through it: events at their fundamental
density, related to each other in the way that makes a flat spacetime. Every region of vacuum contains
a definite number of events, and that number fluctuates as counts do, by about its square root.
\stat{derived}

Such fluctuations of number act on the large-scale geometry like a small cosmological constant whose
size is set by the volume of the observable universe; for the present universe it is of the order of
the observed value~\cite{Ahmed2004}. In Causal Physics this would make dark energy the residue of the
discreteness of the substrate. \stat{open}

The growth of the order on the largest scale is the natural candidate for the expansion of the
universe. \stat{open}

% =============================================================================
\section{Superposition and measurement}\label{sec:onto-superposition}
% =============================================================================

An event can have several admissible continuations, and as long as nothing distinguishes them they
coexist. In Causal Physics there is no separate collapse process; there is stabilisation under
interaction.
\begin{itemize}
\item \emph{Superposition} is the coexistence of admissible causal continuations.
\item \emph{Measurement} is a causal interaction that biases the routing of influence towards one of
  them.
\item \emph{Collapse} is the stabilisation of a single continuation. It is irreversible because
  undoing it would require a coordinated redistribution of influence across an extended region of the
  substrate, and the growing impedance of that region suppresses it. It is a merging that cannot be
  undone.
\end{itemize}
At the discreteness scale, where patterns do not persist long enough to be compressed into
trajectories, every observation perturbs what it probes. Quantum behaviour is then the behaviour of a
substrate that is too unstable for geometric compression, rather than a paradox of measurement.
\stat{interpretation}\quad A quantum dynamics of the order that would make this picture quantitative
is not part of the present work. \stat{open}

% =============================================================================
\section{The frontier}\label{sec:onto-missing}
% =============================================================================

Persistent currents of influence that circulate rather than stream are natural candidates for
\emph{charge} and \emph{spin}: asymmetries in how influence is routed around a structure, invisible from
far away except through their effect on routing. The freedom to relabel routes without changing any
observable flow would then be the origin of \emph{gauge freedom}. These ideas belong to the programme;
none of them is formulated or tested here. \stat{open}

% =============================================================================
\section{The picture in one paragraph}\label{sec:onto-summary}
% =============================================================================

The world is a growing order of events. Its direction is time, its side-by-side structure is space,
and its sameness for every observer is relativity. Influence flows through it and is conserved; its
amount is energy and its direction is momentum. Influence that persists and remembers its direction is
matter, and that memory is inertia; influence in pure transit is light. Matter loads the order, and
the order responds by slowing its clocks, narrowing its cones and adding events so that its throughput
is kept: that response, spreading outwards, is gravity. Where the load is extreme, geometry can no
longer describe the order, and a horizon marks the limit of what the outside can reconstruct. What
cannot be unmerged is entropy, and what has not yet been stabilised is superposition. From events to
gravity each step is tested or derived; self-bound matter, charge, spin and the quantum are the
frontier.
